%Version 3.1 December 2024
% See section 11 of the User Manual for version history
%
%%%%%%%%%%%%%%%%%%%%%%%%%%%%%%%%%%%%%%%%%%%%%%%%%%%%%%%%%%%%%%%%%%%%%%
%%                                                                 %%
%% Please do not use \input{...} to include other tex files.       %%
%% Submit your LaTeX manuscript as one .tex document.              %%
%%                                                                 %%
%% All additional figures and files should be attached             %%
%% separately and not embedded in the \TeX\ document itself.       %%
%%                                                                 %%
%%%%%%%%%%%%%%%%%%%%%%%%%%%%%%%%%%%%%%%%%%%%%%%%%%%%%%%%%%%%%%%%%%%%%

%%\documentclass[referee,sn-basic]{sn-jnl}% referee option is meant for double line spacing

%%=======================================================%%
%% to print line numbers in the margin use lineno option %%
%%=======================================================%%

%%\documentclass[lineno,pdflatex,sn-basic]{sn-jnl}% Basic Springer Nature Reference Style/Chemistry Reference Style

%%=========================================================================================%%
%% the documentclass is set to pdflatex as default. You can delete it if not appropriate.  %%
%%=========================================================================================%%

%%\documentclass[sn-basic]{sn-jnl}% Basic Springer Nature Reference Style/Chemistry Reference Style

%%Note: the following reference styles support Namedate and Numbered referencing. By default the style follows the most common style. To switch between the options you can add or remove �Numbered� in the optional parenthesis. 
%%The option is available for: sn-basic.bst, sn-chicago.bst%  
 
%%\documentclass[pdflatex,sn-nature]{sn-jnl}% Style for submissions to Nature Portfolio journals
%%\documentclass[pdflatex,sn-basic]{sn-jnl}% Basic Springer Nature Reference Style/Chemistry Reference Style
%%\documentclass[pdflatex,sn-mathphys-num]{sn-jnl}% Math and Physical Sciences Numbered Reference Style
%%\documentclass[pdflatex,sn-mathphys-ay]{sn-jnl}% Math and Physical Sciences Author Year Reference Style
%%\documentclass[pdflatex,sn-aps]{sn-jnl}% American Physical Society (APS) Reference Style
\documentclass[pdflatex,sn-vancouver-num]{sn-jnl}% Vancouver Numbered Reference Style — matches Railway Engineering Science's numbered [n] citation style
%%\documentclass[pdflatex,sn-vancouver-ay]{sn-jnl}% Vancouver Author Year Reference Style
%%\documentclass[pdflatex,sn-apa]{sn-jnl}% APA Reference Style
%%\documentclass[pdflatex,sn-chicago]{sn-jnl}% Chicago-based Humanities Reference Style

%%%% Standard Packages
%%<additional latex packages if required can be included here>

\usepackage{graphicx}%
\usepackage{multirow}%
\usepackage{amsmath,amssymb,amsfonts}%
\usepackage{amsthm}%
\usepackage{mathrsfs}%
\usepackage[title]{appendix}%
\usepackage{xcolor}%
\usepackage{textcomp}%
\usepackage{manyfoot}%
\usepackage{booktabs}%
\usepackage{algorithm}%
\usepackage{algorithmicx}%
\usepackage{algpseudocode}%
\usepackage{listings}%
\usepackage{tabularx}%
%%%%

%%%%%=============================================================================%%%%
%%%%  Remarks: This template is provided to aid authors with the preparation
%%%%  of original research articles intended for submission to journals published 
%%%%  by Springer Nature. The guidance has been prepared in partnership with 
%%%%  production teams to conform to Springer Nature technical requirements. 
%%%%  Editorial and presentation requirements differ among journal portfolios and 
%%%%  research disciplines. You may find sections in this template are irrelevant 
%%%%  to your work and are empowered to omit any such section if allowed by the 
%%%%  journal you intend to submit to. The submission guidelines and policies 
%%%%  of the journal take precedence. A detailed User Manual is available in the 
%%%%  template package for technical guidance.
%%%%%=============================================================================%%%%

%% as per the requirement new theorem styles can be included as shown below
\theoremstyle{thmstyleone}%
\newtheorem{theorem}{Theorem}%  meant for continuous numbers
%%\newtheorem{theorem}{Theorem}[section]% meant for sectionwise numbers
%% optional argument [theorem] produces theorem numbering sequence instead of independent numbers for Proposition
\newtheorem{proposition}[theorem]{Proposition}% 
%%\newtheorem{proposition}{Proposition}% to get separate numbers for theorem and proposition etc.

\theoremstyle{thmstyletwo}%
\newtheorem{example}{Example}%
\newtheorem{remark}{Remark}%

\theoremstyle{thmstylethree}%
\newtheorem{definition}{Definition}%

\raggedbottom
%%\unnumbered% uncomment this for unnumbered level heads

\begin{document}

\title[Fair and Transferable Graph RL for Rail Transit Expansion]{Fair and Transferable Graph Reinforcement Learning for Urban Rail Transit Network Expansion Across Asian Cities}

%%=============================================================%%
%% GivenName	-> \fnm{Joergen W.}
%% Particle	-> \spfx{van der} -> surname prefix
%% FamilyName	-> \sur{Ploeg}
%% Suffix	-> \sfx{IV}
%% \author*[1,2]{\fnm{Joergen W.} \spfx{van der} \sur{Ploeg} 
%%  \sfx{IV}}\email{iauthor@gmail.com}
%%=============================================================%%

\author*[1]{\fnm{Quang-Vinh} \sur{Dang}}\email{vinh.dq4@buv.edu.vn}

\affil*[1]{\orgname{British University Vietnam}, \orgaddress{\city{Hung Yen}, \country{Vietnam}}}

%%==================================%%
%% Sample for unstructured abstract %%
%%==================================%%

\abstract{The Metro Network Expansion Problem (MNEP) -- deciding where to extend rail lines and place new stations -- has recently been formulated as a sequential decision process solvable by reinforcement learning (RL). Prior work shows GNN encoders improve demand satisfaction, multi-objective decomposition balances demand against equity and accessibility, and tabular RL matches deep RL at a fraction of the training cost -- but no method combines all three, and none has been tested from small emerging systems to the world's largest rail networks. Emerging systems pose a further difficulty: with only one or two lines, the rail network alone is too sparse to define candidate regions or estimate demand. We propose a GNN-based, multi-objective actor-critic policy for rail transit expansion, scalarizing demand, equity, and accessibility via Tchebycheff decomposition, made sample-efficient via cross-city curriculum transfer and prioritized replay. Each city's bus network is used only as an auxiliary signal, never as an action; the design remains rail-only throughout. We evaluate across Asian cities spanning an order of magnitude in scale: Hanoi, Ho Chi Minh City, Xi'an, and Chengdu. A first, fully trained Xi'an result shows the full method exceeding two architectural ablations on equity and accessibility, and confirms -- after fixing a weight-conditioning bug that had silently masked it -- that one trained policy genuinely varies behavior across trade-offs; on raw demand, a reimplemented tabular RL baseline outperforms ours under a matched budget, consistent with prior efficiency claims at this scale and motivating the larger-scale, cross-city evaluation reported as ongoing work.}

\keywords{metro network expansion, reinforcement learning, graph neural networks, multi-objective optimization, social equity, transfer learning}

%%\pacs[JEL Classification]{D8, H51}

%%\pacs[MSC Classification]{35A01, 65L10, 65L12, 65L20, 65L70}

\maketitle

\section{Introduction}\label{sec1}

A metro line is one of the few pieces of infrastructure a city builds that it can never meaningfully undo. Roads can be repaved, buildings retrofitted, bus routes redrawn overnight -- but once tunnels are bored and viaducts poured, a metro network's shape is fixed for the better part of a century, and every household, employer, and land developer within walking distance of a station reorganizes itself around that fact. Get the sequencing right, and a single well-placed extension can unlock accessibility for a district that spent decades cut off from the rest of the city; get it wrong, and the mistake compounds for generations, entrenching exactly the spatial inequalities that transit investment is supposed to correct. This is why the \emph{Metro Network Expansion Problem} (MNEP) -- given an existing network and a limited construction budget, decide which stations or line segments to add next -- is not merely an interesting combinatorial puzzle. It is one of the highest-stakes recurring decisions in urban planning, and it is made, in most of the world, with far less quantitative rigor than the size of the investment would justify.

Part of the reason is that MNEP is genuinely hard. It is a special case of the broader \emph{Transport Network Design Problem} (TNDP), a family of network-design problems that has occupied operations researchers for decades \citep{farahani2013}, and the metro variant inherits the TNDP's combinatorial explosion: for a city with even a modest number of candidate station sites, the number of ways to sequence an extension budget into a connected line is astronomically large, and evaluating any single candidate network requires simulating how demand redistributes itself across the entire city once that network exists. Classical solution methods -- mixed-integer programming, genetic algorithms, simulated annealing, tabu search, ant colony optimization -- have made real progress on TNDP variants, but nearly all of them purchase tractability by first constraining the search space by hand: restricting candidate stations to a handful of pre-drawn corridors, or fixing the number of lines and their rough geometry before optimization even begins \citep{su2024}. The result is that the algorithm is not really choosing where to build; a planner already did, and the algorithm is polishing the choice.

Reinforcement learning (RL) offers a different bargain. Rather than search a pre-constrained space, an RL agent can be set loose on the full combinatorial problem -- select a station, observe how the network's demand coverage changes, select the next station, repeat -- and learn, through repeated simulated construction, which sequences of decisions tend to produce networks that serve the city well. Three recent papers have each taken a piece of this idea and shown that it works. \citet{su2024} (MetroGNN) represent a city as a graph neural network (GNN) sees it: candidate regions are nodes, spatial adjacency and travel-demand flow are edges, and an attentive actor-critic policy learns to walk this graph, extending or starting lines while a feasibility mask keeps the result physically buildable. The GNN state representation is what makes this work at all -- a policy that can see the city's spatial and demand structure, rather than treating each candidate station as an anonymous slot, substantially outperforms prior heuristics and flat RL baselines at satisfying origin-destination (OD) travel demand. \citet{zhang2024} start from a different frustration: demand-maximizing algorithms have a well-documented tendency to route transit toward wealthy, already-well-connected areas, because that is where the largest raw travel flows are, which is precisely the outcome transit equity policy exists to prevent \citep{hansen1959,raza2023}. They reformulate the reward as a vector -- demand, social equity, and a gravity-style accessibility measure derived from the ``radiation model'' of human mobility \citep{simini2012} -- and combine the three via Tchebycheff scalarization \citep{zhangli2007}, a decomposition technique from multi-objective optimization \citep{hayes2022} that lets a single training run explore genuine trade-offs rather than committing to one fixed weighting up front. Most recently, and most provocatively, \citet{michailidis2026} argue that both of the preceding papers are more machinery than the problem needs: they reformulate MNEP as a Non-Markovian Reward Decision Process that shrinks its effective state-action space, and show that a tabular Monte Carlo RL method -- no neural network, no GPU -- matches deep RL's performance while using an order of magnitude fewer training episodes and a correspondingly smaller carbon footprint. It is a genuinely uncomfortable result for anyone building GNN-based MNEP methods, our own effort included, and it deserves to be taken seriously rather than argued around.

Read together, though, these three papers leave a set of gaps that none of them, individually, was positioned to close. First, no existing method combines the spatial sophistication of a GNN state representation with the fairness-aware, multi-objective reward that \citet{zhang2024} show matters for real-world adoption; a planner today must choose between a policy that understands city structure and one that understands equity, not both at once. Second, \citeauthor{michailidis2026}'s efficiency claim -- that deep RL is unnecessary overhead -- was tested on Xi'an and Amsterdam, both mid-sized, data-rich, already-mature metro systems; whether it still holds at the scale of one of the largest metro networks on Earth, where the state-action space genuinely is enormous, is an open empirical question, not a settled one. Third, and perhaps most consequentially for practice, all three papers train and evaluate a freshly initialized policy for each city they study. This is a reasonable simplification for a first paper on a new method, but it leaves unanswered a question that matters enormously for the cities that most need this technology: does anything learned while training on one city's network transfer to another, or must every city's planning authority independently pay the full training cost from scratch?

That third gap is not an abstract concern for us -- it is a practical blocker. Our motivating case studies are Hanoi and Ho Chi Minh City, Vietnam's two largest metropolitan areas, both mid-expansion right now: Hanoi operates two lines, Ho Chi Minh City one, with several more under construction or in planning. Neither city's existing rail network is remotely large enough, on its own, to define the rich candidate-region structure or estimate the OD demand signal that a GNN-based policy needs to learn from. Prior MNEP methods, developed and validated on cities with a decade or more of accumulated metro infrastructure, simply do not have anything to work with here. We treat this as a data problem rather than a reason to change what the method optimizes: where a city's own rail network is too sparse, we let its existing bus network -- typically far denser, since bus networks are cheaper to build and iterate on -- supply an auxiliary signal for candidate-region definition and demand estimation (Section~\ref{sec:data}). At no point does this let the policy choose, evaluate, or optimize a bus route; the action space, the objective, and the artifact the method ultimately designs remain metro network expansion, full stop. Separately, \citet{kang2019} study how last-train timetables and bus bridging services get coordinated on an \emph{existing} urban railway network -- an operational-phase problem, distinct from the design-phase question we address, but one that reinforces the same point from a different angle: bus and rail networks in practice are never planned in isolation from each other, which is exactly why treating bus data as a legitimate input to rail network design, without letting it become a rail network design \emph{output}, is the right scope boundary to draw.

Putting these threads together, we propose a GNN-based, multi-objective actor-critic policy for urban rail transit network expansion that keeps \citeauthor{su2024}'s spatial representation, keeps \citeauthor{zhang2024}'s fairness-aware multi-objective reward, and directly answers \citeauthor{michailidis2026}'s efficiency critique -- not by abandoning the GNN, but by making the deep-RL training itself cheaper through a cross-city curriculum (pretraining on small networks before fine-tuning on large ones, an idea with no precedent in this literature since no prior MNEP paper trains across more than one city) together with prioritized experience replay \citep{schaul2016} and potential-based reward shaping. We evaluate the resulting method on a set of Asian cities chosen specifically to stress-test the parts of the design space the existing literature has not: Hanoi and Ho Chi Minh City at the small, data-sparse end, and Xi'an and Chengdu -- one of the largest metro networks in the world, and the home city of this journal's publishing institution -- at the large, mature end, with the city set left open-ended rather than fixed, since the value of the transfer-learning claim only grows with a larger and more diverse set of training cities. This paper reports the first result from that programme: a fully implemented, tested, GPU-trained instance of the method on Xi'an, benchmarked directly against a faithful reimplementation of \citeauthor{michailidis2026}'s own tabular method on the same data, together with two ablations that isolate the contribution of the GNN representation and the multi-objective reward respectively. The remaining cities, and the cross-city curriculum itself, are reported as ongoing work in Section~\ref{sec:discussion}, with the infrastructure to complete them already built and described in detail so the path from this paper to that fuller result is concrete rather than aspirational.

In summary, this paper's contributions are:
\begin{enumerate}
\item A GNN-based, multi-objective actor-critic policy for MNEP that, to our knowledge, is the first to combine a heterogeneous spatial-graph state representation with Tchebycheff-decomposed multi-objective reward over demand, equity, and accessibility, and the first to condition the policy itself on the sampled objective-weight vector so a single trained model spans the full trade-off surface rather than committing to one weighting.
\item A cross-city curriculum and sample-efficiency design (prioritized replay, reward shaping) that responds directly to the training-cost and carbon-footprint critique of GNN-based MNEP methods, without abandoning the GNN.
\item A rail-only scope discipline -- using bus network data purely as an auxiliary state signal, never as an action space -- that lets the method extend to small, data-sparse emerging metro systems that no prior MNEP study has evaluated on.
\item A fully real, end-to-end implementation and first experimental result on Xi'an, including two methodological corrections discovered only through actually training the system (Section~\ref{sec:experiments}) that we report explicitly because they generalize beyond this one paper: a naive implementation of weight-conditioned multi-objective RL can silently fail to condition on the weight at all, and naively scalarizing raw per-step reward deltas with a Tchebycheff ideal point is a distinct, easy-to-miss bug from scalarizing the objectives themselves.
\end{enumerate}

\section{Related Work}\label{sec:related}

\subsection{The Transport Network Design Problem}

MNEP inherits its difficulty from the broader Transport Network Design Problem (TNDP), a decades-old family of network-design formulations spanning road, bus, and rail network planning under budget constraints \citep{farahani2013}. \citeauthor{farahani2013}'s review catalogues the dominant solution families -- mixed-integer and bilevel programming, genetic algorithms, simulated annealing, tabu search, ant colony optimization -- and a recurring pattern across nearly all of them: tractability is bought by constraining the search space before optimization begins, typically by restricting candidate infrastructure to a small number of hand-drawn corridors or by fixing the number of lines in advance. This is a reasonable engineering compromise for classical combinatorial optimization, where the size of the true search space (every possible subset and sequencing of candidate stations, evaluated against every possible resulting demand redistribution) is well beyond what exact or even metaheuristic methods can explore directly. It is also precisely the compromise that framing TNDP as a sequential decision process, solvable by reinforcement learning, is designed to remove: an RL agent explores the \emph{full} unconstrained space of station-by-station construction choices, learning through repeated simulated rollouts which sequences tend to produce networks that serve the city well, without a human ever having to pre-select the plausible corridors.

\subsection{Reinforcement learning for MNEP}

Three papers have applied this idea specifically to metro network expansion, and each makes a different methodological bet.

\citet{su2024} (MetroGNN) bet on spatial representation learning. They represent candidate urban regions as a heterogeneous multi-graph -- one edge type for spatial contiguity, one for OD-flow association -- and train an attentive actor-critic policy over this graph, with an action mask enforcing network-design feasibility constraints (station spacing, turn angles, remaining budget) so every rollout produces a physically buildable line. Graph neural networks are, by now, a mature enough tool that their applicability here is unsurprising in retrospect \citep{wu2021gnn}: a candidate region's usefulness as a next station depends on its neighbors' properties as much as its own, which is exactly the kind of relational structure GNNs are built to exploit. The payoff is substantial -- MetroGNN improves OD demand satisfaction over prior heuristic and RL baselines by 15.9\% on average across their evaluated scenarios, exceeding 30\% in higher-budget, more complex cases -- but the method optimizes a single objective (demand satisfaction) and trains a fresh policy from scratch for every city it is applied to.

\citet{zhang2024} bet on fairness instead of representational sophistication. Demand-maximizing policies, theirs included when run in single-objective mode, have a well-known failure mode in transportation planning: because travel demand correlates strongly with existing wealth and connectivity, a policy that purely maximizes satisfied OD flow tends to route new infrastructure toward areas that are already comparatively well served, reproducing and amplifying the very access disparities that public transit investment is nominally meant to correct \citep{hansen1959,raza2023}. \citeauthor{zhang2024} address this by formulating the reward as a three-dimensional vector -- OD demand, a social-equity term, and a gravity-style accessibility measure built on the radiation model of human mobility \citep{simini2012}, which predicts trip flows between locations from population density alone, without requiring a fitted gravity-law exponent -- and combining the three via Tchebycheff decomposition \citep{zhangli2007}, a scalarization technique from the multi-objective optimization literature \citep{hayes2022} whose key property, relative to a simple weighted sum, is that it can reach solutions in non-convex regions of the Pareto front that a weighted sum provably cannot. Sampling a fresh weight vector each training episode, rather than fixing one in advance, lets a single trained policy explore the full range of demand/equity/accessibility trade-offs instead of committing to one planner's a priori judgment about their relative importance. The resulting method outperforms traditional heuristic algorithms (genetic algorithm, ant colony optimization, mathematical programming) by over 30\% and their own weighted-sum RL variant by a smaller margin, evaluated on a real-world Xi'an case study -- but it represents state with hand-engineered features rather than a graph-structured encoder, so it cannot exploit spatial relational structure the way MetroGNN does.

\citet{michailidis2026} bet against both of the preceding papers' machinery. They reformulate MNEP as a Non-Markovian Reward Decision Process, a reformulation that shrinks the effective state-action space enough that a tabular Monte Carlo RL method -- a plain lookup-table policy, no neural network, no GPU -- with fairness-based reward functions matches deep RL performance while using 18$\times$ fewer training episodes and producing an estimated 12$\times$ less CO$_2$ over the training run. This is a pointed and well-evidenced critique of the field's default assumption that more representational capacity is always worth its computational cost, and it is evaluated carefully -- but only on Xi'an and Amsterdam, both mid-sized, data-rich, already-mature metro systems whose state-action spaces, however large in absolute terms, may not be representative of either extreme this paper is concerned with: the enormous scale of the world's largest metro networks, or the acute data sparsity of metro systems still in their first decade of operation.

Read together, these three papers span the design space's most important axes -- spatial representation, multi-objective fairness, computational efficiency -- without any single paper occupying more than one axis at a time, and none reports a result outside a small set of already-mature systems. That gap in \emph{scale coverage}, as much as any single missing feature, is what motivates the case-study selection in Section~\ref{sec:data}.

\subsection{Related operational-phase work on rail-bus coordination}

\citet{kang2019} study a related but structurally different problem: coordinating last-train timetables with bus bridging services on an \emph{existing} urban railway network, via a decomposable mixed-integer linear program, showing that careful coordination of last-train connections across lines together with dispatched bridging buses substantially reduces the number of stranded transfer passengers at service end-of-day. This is an operational-phase problem -- it assumes a fixed rail topology and optimizes how the network is \emph{run}, not what the network \emph{is} -- so we treat it as related work rather than a methodological component of our approach; we neither model bus bridging nor optimize timetables. What it does establish, and what we lean on directly in Section~\ref{sec1}, is that bus and rail service are planned and evaluated jointly in operational practice, not in isolation. That observation is the whole justification for treating a city's existing bus network as a legitimate input signal to rail network \emph{design} (Section~\ref{sec:data}) rather than dismissing it as out of scope for a rail-focused study; \citeauthor{kang2019}'s result is evidence for the premise, not a method we build on.

\section{Problem Formulation}\label{sec:problem}

It is worth being precise about what, exactly, an agent is doing when it ``builds a metro line'' inside this formulation, because the abstraction hides a design decision that matters. The agent does not place stations on a blank map; it walks a graph that already encodes everything we know about the city -- where people live, where they travel to, which regions border which, how well each region is currently served by rail and by bus -- and at each step it commits to exactly one irreversible extension of the network under construction, the same way a real transit authority commits budget to one segment before moving to the next. Framing the problem this way, as a finite-horizon Markov Decision Process (MDP) over a graph of candidate regions, follows \citet{su2024} and \citet{zhang2024} and buys two things classical TNDP solvers do not get for free: the agent can be evaluated and improved through simulated rollouts rather than exhaustive search, and the same formulation, with a different reward vector, is what lets us later condition a single trained policy on a whole family of demand/equity/accessibility trade-offs rather than solving the optimization problem once per trade-off.

Formally, let $\mathcal{N} = \{n_1, \dots, n_k\}$ be the set of candidate regions (Section~\ref{sec:data}). A rail network at step $t$ is a graph $M_t = (V_t, E_t)$ with $V_t \subseteq \mathcal{N}$ the selected stations and $E_t$ the line segments connecting them.

\textbf{State.} $S_t = (M_t, b_t)$: the partial rail network built so far, together with the remaining construction budget. This is deliberately a thin state -- it does not separately track, say, which OD pairs have already been satisfied, because that information is fully recoverable from $M_t$ and the (fixed, precomputed) OD matrix, and keeping the state minimal keeps the GNN encoder's job well-defined: encode the network as it currently stands, not a growing bundle of derived statistics. Bus-derived features (candidate-region richness, demand-proxy signal; Section~\ref{sec:data}) are encoded as node/edge attributes within $M_t$'s graph representation, not as a separate part of the state that could itself be modified by an action -- this is the formal expression of the scope boundary argued for in Section~\ref{sec1}: bus data can inform what the agent sees, never what it is allowed to do.

\textbf{Action.} $A_t \in \mathcal{N} \setminus V_t$: the agent selects one unselected candidate node to add, either extending an existing line or starting a new one. Not every unselected node is a legal choice at every step -- a station that would violate minimum spacing from its neighbors, or force an implausibly sharp turn in the line, is not a station a real transit authority would build -- so the action space is restricted to a feasibility mask $\mathcal{A}(S_t) \subseteq \mathcal{N} \setminus V_t$ enforcing station-spacing and turn-angle constraints and remaining budget, as in \citet{su2024}. $\mathcal{A}(S_t)$ never contains a bus-route action; only rail (metro/tram/light-rail) extensions are ever selectable, per the scope decision in Section~\ref{sec1}.

\textbf{Transition.} Deterministic: $M_{t+1} = M_t \cup \{A_t\}$, with the new station connected to the nearest node on the same line (or starting a new line if none exists within a feasibility-preserving distance). $b_{t+1} = b_t - \text{cost}(A_t)$. An episode terminates when $b_t$ is exhausted or $\mathcal{A}(S_t) = \emptyset$ -- the second condition matters in practice for small, sparse candidate graphs (Hanoi, Ho Chi Minh City) where the agent can paint itself into a corner well before the budget runs out.

\textbf{Reward.} Three network-level objectives are tracked as bounded, cumulative functions of the network built so far: $C^{\text{demand}}(M_t) \in [0,1]$ follows \citeauthor{su2024}'s OD-demand-satisfaction formulation, weighting each OD pair's flow by how directly $M_t$ connects its endpoints; $C^{\text{equity}}(M_t) \in [0,1]$ is one minus a Gini index over per-region accessibility, so that a perfectly equal distribution of access across regions scores 1 and a maximally unequal one scores 0; and $C^{\text{access}}(M_t) \in [0,1]$ is the radiation-model accessibility measure used by \citet{zhang2024}, which predicts the trip flow the network makes available between region pairs from population density alone \citep{simini2012}. All three are computed from the same underlying GNN-encoded region graph (Section~\ref{sec:method}) so they can share intermediate representations rather than requiring three separate forward passes.

\textbf{Scalarization and objective.} Following \citet{zhang2024}, at the start of each training episode a weight vector $\mathbf{w}$ is sampled uniformly from the probability simplex (Algorithm~\ref{alg:curriculum}, line 5), and the three cumulative objectives are combined via augmented Tchebycheff decomposition \citep{zhangli2007},
\begin{equation}
g(\mathbf{C}(M_t) \mid \mathbf{w}, \mathbf{z}^*) = \max_j \left[ w_j \left| C^{(j)}(M_t) - z^*_j \right| \right] + \varepsilon \sum_j w_j \left| C^{(j)}(M_t) - z^*_j \right|,\label{eq:tchebycheff}
\end{equation}
with ideal point $\mathbf{z}^* = (1,1,1)$ and small $\varepsilon > 0$ to exclude weakly-Pareto-optimal solutions, matching the implementation in \texttt{src/rl/reward.py}. Crucially, Tchebycheff's ideal-point comparison is only meaningful against the bounded cumulative objectives $\mathbf{C}(M_t)$ themselves, \emph{not} against their raw per-step deltas -- a per-step demand gain is typically a small fraction of one percent, and comparing it directly to an ideal point of 1 would make every step look equally, uninformatively far from ideal regardless of which action was actually taken. The reward the policy is trained on is therefore the step-to-step change in the already-scalarized, sign-flipped value,
\begin{equation}
r_t = -g(\mathbf{C}(M_{t+1}) \mid \mathbf{w}, \mathbf{z}^*) + g(\mathbf{C}(M_t) \mid \mathbf{w}, \mathbf{z}^*), \label{eq:marginal-reward}
\end{equation}
which telescopes over an episode to the total improvement in scalarized objective value, exactly analogous in spirit to the raw-objective marginal reward $C^{(j)}(M_{t+1}) - C^{(j)}(M_t)$ used by \citet{su2024} for their single demand objective, but applied after scalarization rather than before it. Section~\ref{sec:experiments} reports this distinction not as a footnote but as a real methodological finding: an earlier version of our implementation applied Tchebycheff scalarization to the raw per-objective deltas directly, which trains but silently fails to make the policy responsive to $\mathbf{w}$ at all, since the scalarization term is dominated by the (constant, near-1) distance from a near-zero delta to the ideal point regardless of which objective actually improved.

The policy $\pi_\theta$ is trained to maximize the expected episodic sum of these scalarized rewards, $\mathbb{E}_{\mathbf{w}}\left[\sum_t r_t\right]$, over the sampled weight distribution (undiscounted: each finite-horizon episode ends when the construction budget is exhausted, per the Transition paragraph above). For this expectation to actually depend on $\mathbf{w}$ in a way the policy can exploit, $\mathbf{w}$ must also be visible to the policy itself, not only to the reward function that trains it -- Section~\ref{sec:method-policy} makes this concrete. Done correctly, a single trained policy supports the full range of demand/equity/accessibility trade-offs rather than committing to one fixed weighting at training time; done incorrectly, as our own first implementation attempt was, the weight vector influences the training signal but the resulting policy never learns to look at it, and behaves identically regardless of which trade-off is requested at deployment time. We return to exactly how we detected and fixed this in Section~\ref{sec:experiments}, because we think the failure mode is general enough to be worth naming for other implementers of weight-conditioned multi-objective RL, not specific to our codebase.

\section{Method}\label{sec:method}

The method has three pieces, each answering a different one of the gaps identified in Section~\ref{sec1}. A graph neural network encoder (Section~\ref{sec:method-gnn}) gives the policy access to the same spatially-structured city representation that makes MetroGNN effective, extended to also see bus-informed auxiliary signal where a city's rail network alone is too sparse. An attentive, weight-conditioned actor-critic policy (Section~\ref{sec:method-policy}) turns that representation into station-by-station construction decisions, and is explicitly built so that the Tchebycheff weight vector sampled each episode (Section~\ref{sec:problem}) is visible to the policy itself, not only to the reward that trains it -- the distinction that Section~\ref{sec:problem} flagged as a real, easy-to-miss implementation pitfall. A cross-city curriculum with prioritized replay and reward shaping (Section~\ref{sec:method-curriculum}) is the direct response to \citeauthor{michailidis2026}'s efficiency critique: rather than accept that GNN-based MNEP methods must be expensive to train, we ask whether training cost can be amortized across cities the way it already is, informally, across the training episodes within one city.

\subsection{GNN state encoder}\label{sec:method-gnn}

We encode $M_t$ with a two-relation heterogeneous graph neural network, following \citeauthor{su2024}'s scheme (\citeyear{su2024}, Sec.~3.1--3.2), implemented in \texttt{src/gnn/encoder.py}. Each candidate region $n_i$ has a raw feature vector $\mathbf{a}_i$ (population density, income/employment where available, bus-derived candidate-region and demand-proxy signal from Section~\ref{sec:data}, and existing-rail proximity), projected to an initial hidden state $\mathbf{h}_i^{(0)} = \mathbf{W}_A \mathbf{a}_i$. Two edge types connect candidate regions: spatial-contiguity edges $e^s_{ij}$ (shared boundary, Section~\ref{sec:data}'s queen-contiguity construction) and OD-flow edges $e^o_{ij}$ (disaggregated travel demand). At each layer $l$, messages are aggregated separately per relation and combined with a residual connection:
\begin{align}
\mathbf{h}_{s,i}^{(l)} &= \sum_{j} e^s_{ij} \, \mathbf{W}_s^{(l)} \mathbf{h}_j^{(l)}, \qquad
\mathbf{h}_{o,i}^{(l)} = \sum_{j} e^o_{ij} \, \mathbf{W}_o^{(l)} \mathbf{h}_j^{(l)}, \label{eq:gnn-relations}\\
\mathbf{h}_i^{(l+1)} &= \tanh\!\left( \mathbf{W}_c^{(l)} \left[\mathbf{h}_{s,i}^{(l)} \,\|\, \mathbf{h}_{o,i}^{(l)}\right] + \mathbf{h}_i^{(l)} \right), \label{eq:gnn-combine}
\end{align}
where $\|$ denotes concatenation. Stacking $L$ such layers yields final node embeddings $\mathbf{h}_i^{(L)}$ that jointly encode a region's rail/OD structure and its bus-informed auxiliary features -- the bus signal enters only as an entry of $\mathbf{a}_i$, never as its own node or edge type, so it cannot be selected as an action regardless of how the graph is encoded.

\subsection{Multi-objective policy}\label{sec:method-policy}

The actor-critic policy (\texttt{src/rl/policy.py}) operates on the embedded graph $\{\mathbf{h}_i^{(L)}\}$. A query vector $\mathbf{q}_t$ represents the current partial network, computed by mean-pooling the embeddings of already-selected nodes $V_t$ (or all candidate nodes if $V_t = \emptyset$, at the start of an episode). Each candidate node is scored by scaled dot-product attention against $\mathbf{q}_t$,
\begin{equation}
\text{score}(n_i) = \frac{(\mathbf{W}_k \mathbf{h}_i^{(L)}) \cdot (\mathbf{W}_q \mathbf{q}_t)}{\sqrt{d}}, \label{eq:policy-score}
\end{equation}
with $d$ the embedding dimension, and infeasible actions ($n_i \notin \mathcal{A}(S_t)$) masked to $-\infty$ before a softmax yields the action distribution $\pi_\theta(A_t \mid S_t)$. A separate value head, a two-layer MLP over $\mathbf{q}_t$, estimates $V_\phi(S_t)$ for the actor-critic baseline -- the general actor-critic template we follow, learning a policy and a value function jointly rather than either alone, goes back to \citet{konda2000}, and remains the standard choice whenever the action space is large enough that pure value-based methods (learning a value for every state-action pair directly) become impractical, as it is here.

Conditioning the policy on the sampled Tchebycheff weight vector $\mathbf{w}$ is what makes the multi-objective claim in Section~\ref{sec:problem} actually true rather than aspirational, and it is worth stating precisely because it is easy to get subtly wrong: it is not enough for $\mathbf{w}$ to shape the \emph{reward} the policy is trained on, because a policy trained on a randomly varying reward signal, with no way to observe which variant it is currently being asked to optimize, can at best learn some compromise behavior averaged over the sampled weight distribution -- it has no mechanism to specialize its behavior to any particular $\mathbf{w}$, since $\mathbf{w}$ never enters its input. We embed $\mathbf{w}$ through a small learned linear layer and add it directly to the pooled query, $\mathbf{q}_t' = \mathbf{q}_t + \mathbf{W}_e \mathbf{w}$, before it is used both for attention scoring (Eq.~\ref{eq:policy-score}, with $\mathbf{q}_t'$ in place of $\mathbf{q}_t$) and for the value estimate. This is a small architectural change -- one additional linear layer -- but Section~\ref{sec:experiments} shows it is the difference between a policy that behaves identically under demand-only, equity-only, and coverage-only weightings and one that genuinely does not.

\subsection{Cross-city curriculum and sample efficiency}\label{sec:method-curriculum}

\citeauthor{michailidis2026}'s central empirical claim -- that a tabular RL method, needing no neural network and no GPU, matches deep RL's performance at a fraction of the training cost and carbon footprint -- is not a claim we set out to refute. It is a claim we take as a design constraint: if the value of a GNN-based, multi-objective policy is going to justify its additional computational cost, that value has to come from somewhere real, not from ignoring an efficiency argument that the evidence otherwise supports. Two places suggest themselves. The first is scale: a lookup table's efficiency advantage on a small, mid-sized network like Xi'an's may not hold once the state-action space is as large as Chengdu's, where the sheer number of candidate regions could make tabular methods' implicit assumption -- that most state-action pairs will be visited often enough to estimate their value reliably -- start to break down, though establishing this is squarely future work (Section~\ref{sec:discussion}) rather than a claim this paper's Xi'an-only result can support. The second, and the one we act on directly here, is that a GNN encoder's training cost need not be paid independently by every city, the way both \citeauthor{michailidis2026}'s tabular method and every one of the three reference RL papers effectively require: a lookup table learned on one city's grid has no structure that transfers to a different city's grid, but a GNN's learned message-passing weights are, in principle, city-agnostic, since they operate on the same relative notions of spatial contiguity and OD-flow structure regardless of which city's graph they are applied to.

We therefore introduce a cross-city curriculum: policy weights are pretrained on the smaller networks (in our full four-city design, Hanoi and Xi'an) and used to initialize training on the larger networks (Ho Chi Minh City, Chengdu), rather than training each city's policy from scratch, following the general intuition -- well established in supervised deep learning, less commonly exploited across environments in RL -- that a network which has already learned useful low-level structure converges faster from that starting point than from a random initialization. This is combined with prioritized experience replay \citep{schaul2016}, which samples training transitions in proportion to the magnitude of their temporal-difference error rather than uniformly, so the limited training budget is spent disproportionately on the transitions the policy is currently getting most wrong, and potential-based reward shaping, which accelerates credit assignment for the delayed, sparse parts of the reward signal without altering the optimal policy the agent ultimately converges to. Algorithm~\ref{alg:curriculum} states the full procedure; Section~\ref{sec:experiments} reports results from the first city in this sequence, Xi'an, trained from a random initialization since no smaller city has yet been trained to transfer from -- the curriculum step itself, and the efficiency gains it is expected to produce, remain to be measured once a second city's training run is complete (Section~\ref{sec:discussion}).

\begin{algorithm}
\caption{Cross-city curriculum training for the GNN multi-objective policy}\label{alg:curriculum}
\begin{algorithmic}[1]
\Require Cities ordered by network scale $C_1 \le C_2 \le \dots \le C_n$; Tchebycheff weight set $W$
\State Initialize policy/value network parameters $\theta_0$ randomly
\For{$i = 1, \dots, n$}
    \State Initialize $\theta \gets \theta_{i-1}$ \Comment{transfer weights from the previous, smaller city}
    \For{each training episode on city $C_i$}
        \State Sample Tchebycheff weight vector $w \sim W$
        \State Roll out policy on $C_i$'s region graph, collect transitions $(s_t, a_t, r_t, s_{t+1})$
        \State Store transitions in prioritized replay buffer $\mathcal{D}_i$
        \State Sample prioritized batch from $\mathcal{D}_i$, apply potential-based reward shaping
        \State Update $\theta$ via actor-critic gradient step on the scalarized reward
    \EndFor
    \State $\theta_i \gets \theta$
\EndFor
\State \Return $\theta_n$
\end{algorithmic}
\end{algorithm}

\subsection{Instantiating the action space per city}\label{sec:method-implementation}

Section~\ref{sec:problem} defines the action space abstractly, as selecting any feasible unselected candidate node $A_t \in \mathcal{A}(S_t)$. In practice, how that feasibility mask is realized depends on how a given city's candidate regions are represented, and it is worth being explicit that this varies across our case-study cities rather than presenting a single uniform implementation as if it were forced by the formulation. For Xi'an (Section~\ref{sec:experiments}), we adopt the standard benchmark environment for this exact city and dataset -- an 8-directional grid walk, where the agent occupies one grid cell and may extend the line to any of its eight compass neighbors subject to the feasibility mask -- because this is the action-space convention used by \citet{michailidis2026}'s own public environment, and matching it exactly is what makes a direct, apples-to-apples comparison against their reported numbers possible rather than merely approximate. For Ho Chi Minh City's commune/ward-level candidate regions (Section~\ref{sec:data}), where regions are irregular polygons rather than uniform grid cells, the natural realization is instead \citeauthor{su2024}'s formulation: the action space is any spatially-contiguous unselected neighbor of the current line's endpoint, which for an irregular polygon graph is exactly analogous to what an 8-directional grid walk is for a regular one. Both are legal specializations of the same abstract MDP in Section~\ref{sec:problem}; neither changes the reward, the objectives, or the GNN encoder's architecture. We flag this explicitly because it is the kind of implementation detail that is easy to leave implicit in a way that makes results across cities look more directly comparable than they actually are -- Section~\ref{sec:discussion} returns to what this does and does not license us to conclude when comparing across cities with different action-space realizations.

\section{Case Studies and Data}\label{sec:data}

The four cities anchoring this evaluation were not chosen for convenience; each one is doing a specific piece of argumentative work that the reference literature's city choices could not. Xi'an, a provincial capital of roughly thirteen million people in Shaanxi, China, has already been studied by two of our three reference papers -- \citet{michailidis2026} used it to argue tabular RL suffices, and \citet{zhang2024} used it to demonstrate the value of multi-objective fairness -- which makes it, unusually for this literature, a city where our own results can be checked against two independently reported baselines on the same underlying network rather than only our own reimplementations of them. Chengdu, roughly the same population but with one of the most extensive metro networks anywhere in the world after a decade of construction that shows no sign of slowing, is the scale extreme neither prior paper tested: if a tabular method's efficiency advantage is going to break down anywhere, or if a GNN's spatial reasoning is going to earn its computational cost anywhere, Chengdu's state-action space is large enough that the answer should become visible. Chengdu is also, not incidentally, the home city of Southwest Jiaotong University, this journal's publishing institution, and a network its own transportation researchers know intimately -- a reason to hold our Chengdu results to a higher standard of scrutiny, not a lower one. Hanoi and Ho Chi Minh City are the opposite extreme, and the one the reference literature has never touched at all: Vietnam's two largest metropolitan areas, home to a combined seventeen million people, are only now building the metro networks that Xi'an and Chengdu have had for a decade or more, which makes them exactly the kind of small, data-sparse, actively-expanding system this paper's motivating scope decision (Section~\ref{sec1}) was designed for -- and exactly the kind of system where getting network expansion planning right, before path-dependence sets in, matters most.

Because MNEP RL methods are trained per city rather than pooled across a fixed benchmark set, the city set is deliberately not fixed at four architecturally -- nothing in the method requires exactly this set, and we extend it with additional Asian cities (candidates: Beijing and Changsha, the two cities used by \citet{su2024}, if their released code/data prove usable) wherever open data can be confirmed, since a larger and more diverse city set strengthens both the transfer-learning claim of Section~\ref{sec:method-curriculum} and the scale-dependent generalization claim motivating the Chengdu comparison. Table~\ref{tab:cities} summarizes each city's role and its data-readiness status as of this paper's first experimental result; Section~\ref{sec:experiments} reports results for Xi'an, the furthest along, with the remaining three at varying stages of the pipeline described below.

\begin{table}[h]
\caption{Case study cities, their role in the experimental design, and confirmed open-data status (see text for sources)}\label{tab:cities}
\begin{tabularx}{\textwidth}{@{}l X X X@{}}
\toprule
City & Role & Scale (approx.) & Data status \\
\midrule
Hanoi & Novel, small/emerging metro & 2 lines + planned extensions & Network/boundaries strong; OD is proxy-only \\
Ho Chi Minh City & Novel, medium/growing metro & 1 line (2024) + planned lines & Strongest of the four: real shapefiles for metro, bus, 2018 OD survey, population, POI, census (confirmed) \\
Xi'an & Direct benchmark vs.\ \citet{michailidis2026} and \citet{zhang2024} (both use Xi'an) & Medium, previously studied & OD grid confirmed reusable from prior code release \\
Chengdu & Scale stress-test & $\sim$500+ km, one of the world's largest & Network geometry confirmed (multi-year); ridership portal unconfirmed \\
\botrule
\end{tabularx}
\end{table}

\subsection{Data sources per city}

The sources below have been directly verified by inspecting the underlying repositories (not merely located), except where noted.

\textbf{Network topology and planned extensions.} Existing rail geometry for all four cities is obtainable from OpenStreetMap/Overpass. For Xi'an and Chengdu specifically, a public dataset covering 38 Chinese cities provides \emph{year-by-year} (2010--2020) station-point and line shapefiles, confirmed by direct inspection; because it is multi-year, it additionally enables a validation protocol not used by any of the three reference RL papers -- checking whether the policy's highest-priority predicted extensions, given an early-year network as the starting state, match what was actually built by a later year. Planned-extension maps exist as official planning diagrams for all four cities (Hanoi's 8-line master plan; Ho Chi Minh City's MAUR master plan; Xi'an's and Chengdu's Five-Year-Plan transport documents) but require manual digitization -- none is available as a ready-to-use shapefile.

\textbf{OD demand.} Ho Chi Minh City has by far the strongest confirmed source: the public \texttt{CityScope/CSL\_HCMC} repository (MIT Media Lab) ships real shapefiles -- metro routes/stations, bus routes/stops, a 2018 household-travel-survey OD matrix, district- and ward-level population, points of interest (POI), and ward-level census with age breakdown -- verified by direct inspection, though the repository carries no explicit license, so we treat it as research/fair-use only and do not redistribute the underlying files. For Xi'an, a public MIT-licensed environment (via \citet{michailidis2026}'s code release and its \texttt{mo-tndp} submodule) ships a ready-to-use 29$\times$29-cell grid with 463{,}010 nonzero OD pairs, the existing metro lines as grid-coordinate paths, and per-cell house-price/equity groupings -- confirmed by direct inspection. That dataset's own citation file attributes it to \citet{michailidis2024} rather than to the KDD~2020 paper we had provisionally assumed as its origin; we cite the confirmed source and do not assert the earlier, unverified lineage. Notably, \citet{zhang2024} -- one of our three reference RL baselines -- also uses Xi'an as its case study, so our Xi'an results can be compared against two, not just one, of the three reference methods' originally-reported numbers. For Hanoi and Chengdu (and as a fallback for the others), we construct a gravity-model demand proxy from WorldPop population-density rasters and OpenStreetMap/Amap/Baidu POI density, calibrated against published ridership figures (e.g., Hanoi's Cat Linh--Ha Dong line: 13.68M riders in 2025, up 15.3\% year-on-year; Ho Chi Minh City Line 1: 20.56M riders in 2025).

\textbf{Bus network (auxiliary signal only).} Each city's existing bus route network is used solely as an auxiliary input to the rail expansion problem: it helps define a richer set of candidate expansion regions in cities where the rail network alone is too sparse (Hanoi, Ho Chi Minh City) and contributes an additional demand-proxy feature. Bus routes are never an output or design target of the method. Ho Chi Minh City's bus network is available as a real shapefile (CSL\_HCMC, above); the remaining cities use OpenStreetMap bus-route extraction. \emph{Coverage caveat, confirmed by direct inspection:} CSL\_HCMC's 189 bus stops fall entirely within a small downtown/pilot-study area, giving nonzero bus-derived signal for only $\sim$32 of the city's 322 candidate zones ($\sim$10\%); outside that area the bus feature is uniformly zero, and the OpenStreetMap fallback (or a wider bus dataset, if found) will be needed for full-city coverage.

\textbf{Administrative boundaries / candidate regions.} GADM (the Database of Global Administrative Areas) provides consistent district-level administrative boundaries for all four cities, usable to define MetroGNN-style candidate expansion regions.

\textbf{Coordinate systems.} Chinese POI sources (Amap, Baidu) report coordinates in the obfuscated GCJ-02 system, which must be converted to WGS84 before merging with OSM/metro station coordinates; this conversion step is required for Xi'an and Chengdu but not for the OSM-only Vietnamese cities.

%% TODO: inspect github.com/tsinghua-fib-lab/MetroGNN (Beijing/Changsha data) as a candidate
%% 5th/6th city for the open-ended city set.
%% TODO: manually check data.chengdu.gov.cn for a metro passenger-flow open-data catalog entry.
%% Full source-by-source detail: docs/data_sources.md

\section{Experiments and Results}\label{sec:experiments}

We report results from the first city in the evaluation programme described in Section~\ref{sec:data}: Xi'an, chosen first because it is the one city with a ready-to-use, previously-published OD environment (Section~\ref{sec:data}), which let us validate the full pipeline -- data loading, GNN encoding, multi-objective policy training, and baseline comparison -- against a real, external, previously-studied dataset before extending it to cities where we ourselves had to assemble the data pipeline from raw sources. What follows is accordingly a single-city result, honestly scoped as such throughout; the cross-city curriculum described in Section~\ref{sec:method-curriculum}, which is where we expect our method's efficiency advantage over training each city from scratch to actually materialize, has not yet been measured, since it requires a second trained city to transfer from. Section~\ref{sec:discussion} returns to exactly what this first result does and does not license us to conclude.

\subsection{Baselines}

We compare against three baselines. The tabular RL baseline is a direct, literal execution of \citeauthor{michailidis2026}'s own public code release (Section~\ref{sec:data}) on the same Xi'an data our method uses, not a reimplementation from the paper's description -- this is the strongest form of baseline comparison available, since any performance difference cannot be attributed to reimplementation drift. The other two reference methods, MetroGNN \citep{su2024} and the multi-objective actor-critic of \citet{zhang2024}, have no public code release we could locate and run directly, so we use the two natural ablations of our own method as principled stand-ins rather than either omitting the comparison or reimplementing from scratch under time pressure that would itself risk introducing errors: \texttt{single\_objective} (our full GNN architecture, trained with the Tchebycheff weight fixed to $(1,0,0)$ every episode, i.e.\ demand only) approximates MetroGNN's single-objective, GNN-based design; \texttt{flat\_encoder} (our full multi-objective reward and training procedure, but with the GNN encoder replaced by a linear projection with no message passing) approximates \citeauthor{zhang2024}'s multi-objective, non-graph-structured design. We are explicit that these are principled approximations of the two papers' \emph{design choices}, not reproductions of their \emph{reported numbers} -- our environment, feature set, and training budget all differ from theirs, so a numerical difference between our \texttt{single\_objective} row and \citeauthor{su2024}'s published results reflects those differences as much as anything about single-objective GNN policies in general. What the two ablations \emph{do} let us isolate cleanly, because everything except one architectural choice is held fixed, is the marginal contribution of the GNN representation and the marginal contribution of the multi-objective reward within our own method.

\subsection{Metrics}

OD demand satisfied (\%), read directly from the environment's own reward signal (Section~\ref{sec:method-implementation}); social equity, computed as one minus the Gini index over the five price-based population groups' cumulative satisfied demand (Section~\ref{sec:problem}); coverage, the fraction of grid cells within Chebyshev distance $K=2$ of any built station. We flag explicitly, and repeat in Section~\ref{sec:discussion}'s limitations, that coverage is a simplified accessibility \emph{proxy} we introduce for this experiment, not \citeauthor{zhang2024}'s radiation-model accessibility measure \citep{simini2012} -- implementing the full radiation model is a well-scoped, tractable piece of near-term follow-up work, not a fundamental barrier, but reporting a proxy as though it were the exact published measure would misrepresent what was actually computed, so we do not. Training episodes and wall-clock GPU time are reported directly; a formal CO$_2$ accounting following \citeauthor{michailidis2026}'s method is deferred to the full multi-city evaluation (Section~\ref{sec:discussion}), since a meaningful carbon comparison needs the cross-city curriculum's amortized cost, not one city's isolated training run.

\subsection{Experimental setup}

All four rows of Table~\ref{tab:results} are trained under matched conditions: 500 training episodes, 20 stations per episode, on the real Xi'an 29$\times$29 candidate grid with existing metro lines included (Section~\ref{sec:data}). The tabular baseline uses its own paper's default hyperparameters; our method and both ablations share the hyperparameters in Appendix~\ref{secA1}. Full reproduction commands, including a real upstream bug we found and patched in the tabular baseline's own evaluation code, are recorded in \texttt{docs/baseline\_reproduction.md} in the project repository, described in detail below because we think the two bugs we found while getting this experiment running are informative beyond this one paper.

\textbf{Bug 1: a state-tracking error in the baseline's own test-time evaluation.} \citeauthor{michailidis2026}'s public \texttt{QLearningTNDP.test()} method took the agent's current grid position from the raw environment observation returned by \texttt{env.reset()}/\texttt{env.step()} (a \texttt{MultiBinary} vector encoding covered cells) rather than from \texttt{info['location\_grid\_index']} (a scalar grid index), as their own \texttt{train()} method -- two methods above it in the same file -- correctly does. A commented-out line in \texttt{test()} reveals that the intended fix was already known but never applied. The bug caused an \texttt{IndexError} on every attempted evaluation run; the training loop itself was unaffected and required no changes. We patched \texttt{test()} to match \texttt{train()}'s own convention and verified the fix reproduces sensible, converged policies (Section~\ref{sec:experiments} results below). We report this not to criticize a paper we otherwise rely on and whose central efficiency claim we take seriously, but because it is exactly the kind of small, easy-to-miss discrepancy between a paper's training code and its evaluation code that silently invalidates published benchmark numbers if it goes unnoticed, and because open, reusable research code -- which \citeauthor{michailidis2026} did the right thing by publishing -- is only as valuable as its verifiability.

\textbf{Bug 2: our own weight-conditioning gap.} Our first working version of the training loop reproduced the design in Section~\ref{sec:problem} on paper but not in code: the sampled Tchebycheff weight vector $\mathbf{w}$ shaped the training reward exactly as intended, but was never actually passed to the policy network, so the policy had no way to observe which trade-off it was being trained toward at any given moment. We caught this only because we built a direct sanity check into the training script: after training, evaluate the same trained model greedily under four different weight vectors (uniform, demand-only, equity-only, coverage-only) and check whether the resulting constructed lines actually differ. They did not -- byte-for-byte identical demand, equity, and coverage scores and an identical action sequence across all four weightings, which is conclusive evidence that the ``single policy, many trade-offs'' claim was false for that implementation, despite the training loss decreasing normally and the training-time metrics looking entirely unremarkable. Section~\ref{sec:method-policy}'s query-conditioning fix ($\mathbf{q}_t' = \mathbf{q}_t + \mathbf{W}_e\mathbf{w}$) resolves it; Table~\ref{tab:weight-sensitivity} below reports the same four-weighting sanity check on the corrected model, where the four settings now produce genuinely different lines. We think this failure mode -- a multi-objective RL method that trains without error, reports plausible-looking metrics, and is nonetheless not actually multi-objective in the sense claimed -- is worth naming explicitly for other implementers of weight-conditioned policies, since nothing about the failure is visible from the training curves alone; only an explicit behavioral check across weight settings reveals it. We would not have caught it, and would have reported a Table~\ref{tab:results} row for a method that silently was not doing what the paper claims, without deliberately building in this check.

\subsection{Ablations}

Three comparisons isolate the contribution of each design choice, all trained and evaluated under the setup above: \texttt{full} vs.\ \texttt{single\_objective} isolates the value of the multi-objective reward (holding the GNN architecture fixed); \texttt{full} vs.\ \texttt{flat\_encoder} isolates the value of the GNN's spatial message passing (holding the multi-objective reward fixed); and \texttt{full} vs.\ the tabular RL baseline is the direct test of \citeauthor{michailidis2026}'s efficiency argument under matched training budget on the one axis their own paper measures, raw demand satisfaction.

\begin{table}[h]
\caption{Xi'an results, greedy evaluation under uniform (1/3, 1/3, 1/3) Tchebycheff weights, 500 training episodes, 20 stations. Tabular RL: mean $\pm$ sample std over 5 seeds (1, 2, 3, 7, 42). Our three variants: mean $\pm$ sample std over 3 seeds (1, 2, 42) each -- fewer seeds than the tabular baseline, disclosed as a limitation below. Demand for our method's rows is deliberately not directly comparable to the tabular baseline's, since ours is trained to jointly optimize three objectives under a randomly varying weight each episode, while the tabular baseline is trained purely to maximize demand -- see discussion below.}\label{tab:results}
\begin{tabularx}{\textwidth}{@{}X ccc c@{}}
\toprule
Method & Demand (\%) & Equity & Coverage & Training episodes \\
\midrule
Tabular RL \citep{michailidis2026} & $1.36 \pm 0.56$ & n/a$^\dagger$ & n/a$^\dagger$ & 500 \\
\texttt{single\_objective} (MetroGNN-like) & $0.47 \pm 0.46$ & $0.45 \pm 0.11$ & $0.14 \pm 0.01$ & 500 \\
\texttt{flat\_encoder} (Zhang et al.-like) & $0.24 \pm 0.25$ & $0.29 \pm 0.08$ & $0.11 \pm 0.05$ & 500 \\
\texttt{full} (Ours) & $0.40 \pm 0.37$ & $0.55 \pm 0.17$ & $0.15 \pm 0.03$ & 500 \\
\botrule
\end{tabularx}

\vspace{0.5em}
\footnotesize $\dagger$ The tabular baseline's own reward formulation optimizes demand only (\texttt{max\_efficiency}); it does not compute an equity or coverage metric in units comparable to ours, so we report those cells as not applicable rather than compute a number that would misrepresent what the baseline actually optimizes.
\end{table}

\begin{table}[h]
\caption{Weight-conditioning sanity check: the same trained \texttt{full} model (seed 42), evaluated greedily under four different Tchebycheff weight vectors. Non-identical rows confirm the policy genuinely conditions on $\mathbf{w}$ (Section~\ref{sec:method-policy}), unlike our first, buggy implementation.}\label{tab:weight-sensitivity}
\begin{tabular}{@{}lccc@{}}
\toprule
Weight vector $\mathbf{w}$ & Demand (\%) & Equity & Coverage \\
\midrule
Uniform (1/3, 1/3, 1/3) & 0.39 & 0.536 & 0.186 \\
Demand-only (1, 0, 0) & 0.39 & 0.536 & 0.186 \\
Equity-only (0, 1, 0) & 0.43 & 0.511 & 0.167 \\
Coverage-only (0, 0, 1) & 0.28 & 0.399 & 0.120 \\
\botrule
\end{tabular}
\end{table}

Table~\ref{tab:results} reflects the complete picture across three seeds per method rather than the single-seed snapshot we first drafted this section around, and we report the fuller picture in full, including where it complicates a cleaner single-seed story. Four findings stand out, all reported plainly including where they cut against the paper's own hypotheses, because we would rather a reader trust the honesty of a small preliminary result than doubt the completeness of a larger one.

First, and most immediately strikingly, the seed-to-seed variance is large -- large enough that it changes which conclusions the data actually supports. Demand for \texttt{single\_objective} ranges from 0.10\% to 0.98\% across three seeds; for \texttt{full}, from 0.03\% to 0.78\%. This is a genuine property of training a several-thousand-parameter policy with vanilla REINFORCE for only 500 episodes on a small candidate graph, not an artifact of our reporting, and it is the single most important limitation of this result (elaborated below): with three seeds, none of the demand differences between \texttt{full}, \texttt{single\_objective}, and \texttt{flat\_encoder} in Table~\ref{tab:results} are large relative to their standard deviations, and we do not claim otherwise.

Second, on the one metric directly comparable in spirit across all four rows -- how much OD demand a 20-station, 500-episode-trained line satisfies -- the tabular baseline's mean (1.36\%) comfortably exceeds all three of our variants' means (0.24--0.47\%), by a factor of roughly three to six depending on which of our variants it is compared against. This is not the result we expected going in, and we do not think it should be read as evidence against a GNN-based, multi-objective approach in general; it is, we think, mostly explained by two features of this specific comparison that will not hold once the method is evaluated as designed. Our variants spend their 500 episodes learning to balance either three objectives under a randomly sampled weight each episode (\texttt{full}, \texttt{flat\_encoder}) or a single objective but through a several-thousand-parameter neural policy rather than a lookup table (\texttt{single\_objective}), while a tabular method with one entry per state-action pair can, on a state-action space as constrained as a single 20-station Xi'an rollout, simply memorize good trajectories faster than any of our neural variants can learn a generalizable representation from the same number of samples -- exactly \citeauthor{michailidis2026}'s point, and exactly the pattern we would expect to reverse once the state-action space is large enough (Chengdu) or the network is initialized from a useful prior (the cross-city curriculum) rather than trained from scratch on one small city with a high-variance policy-gradient method.

Third, on equity and coverage -- the two objectives the tabular baseline's own reward formulation does not optimize for at all -- \texttt{full}'s mean exceeds both ablations': $0.55 \pm 0.17$ equity and $0.15 \pm 0.03$ coverage, versus \texttt{single\_objective}'s $0.45 \pm 0.11$ and $0.14 \pm 0.01$ and \texttt{flat\_encoder}'s $0.29 \pm 0.08$ and $0.11 \pm 0.05$. The equity gap between \texttt{full} and \texttt{flat\_encoder} in particular ($0.55$ vs.\ $0.29$, roughly two combined standard deviations apart) is the single most suggestive finding in this table: it is consistent with the hypothesis that a graph-structured state representation genuinely helps a multi-objective policy balance fairness against demand, not merely that more parameters help in general, since \texttt{single\_objective} has the same parameter count as \texttt{full} and does not show the same equity advantage over \texttt{flat\_encoder}. We call this suggestive rather than established, deliberately, given three seeds and one city.

Fourth, and most methodologically important independent of any of the above, Table~\ref{tab:weight-sensitivity} confirms the corrected weight-conditioning mechanism works as designed -- the four weightings produce genuinely different lines and scores, not the identical output the buggy first version produced -- but it also surfaces a counterintuitive result worth reporting rather than smoothing over: coverage-weighted evaluation produces the \emph{lowest} coverage score of the four settings (0.120, versus 0.186 under uniform weights), which on inspection traces to the specific action sequence the coverage-only policy converged to -- a long, straight run in a single compass direction, which packs stations close together and therefore covers a narrower corridor than a policy that spreads stations out, exactly the kind of locally-plausible-but-globally-suboptimal behavior a policy trained for only 500 episodes under a rarely-sampled weight setting (coverage-only weighting is one of infinitely many points on the simplex, most training episodes sample something else) has had comparatively little experience actually optimizing. We read this as an honest signature of under-training on a rarely-visited part of weight space, not as a flaw in the weight-conditioning mechanism itself -- the mechanism visibly works (Table~\ref{tab:weight-sensitivity}'s four rows are not identical, which is the property we set out to verify), even though what it has learned to do under that particular weighting is not yet good.

\section{Discussion}\label{sec:discussion}

This paper set out to close three gaps in the MNEP literature -- no method combining GNN spatial representation with multi-objective fairness, no test of whether tabular RL's efficiency advantage holds at scale extremes, and no test of whether learned expansion policies transfer across cities -- and it is important to be precise about which of the three this first, single-city, single-seed Xi'an result actually speaks to, and which remain open.

\textbf{On the first gap}, the result is genuinely informative, if modest and mixed rather than uniformly favorable: Table~\ref{tab:results} shows \texttt{full} exceeding both ablations' means on equity and coverage, with the equity gap over \texttt{flat\_encoder} the most suggestive single finding in this paper, but demand itself is statistically indistinguishable across all three variants given three-seed variance this high -- an honest result, not the clean sweep a single seed had suggested before we ran the additional seeds. Table~\ref{tab:weight-sensitivity} demonstrates, independent of that variance and after fixing a real bug that would otherwise have made this claim false, that a single trained policy meaningfully varies its behavior across the demand/equity/coverage trade-off space rather than collapsing to one fixed compromise. Neither finding required Chengdu-scale data or a multi-city curriculum to establish; both are properties of the method's architecture and training procedure that a single well-instrumented city can reveal -- though the demand finding specifically shows why more than one seed is needed before treating any single-run result as conclusive.

\textbf{On the second gap} -- whether tabular RL's efficiency advantage holds at scale -- this Xi'an result, if anything, is a data point \emph{for} \citeauthor{michailidis2026}'s position rather than against it: under a matched, modest training budget on a mid-sized city, the tabular method's demand score comfortably beats ours. We think this is the honest and expected outcome at Xi'an's scale, and we say so plainly rather than searching for a more favorable metric to headline, because the paper's actual contribution to this question is a hypothesis worth testing, not a result yet: that the tabular method's advantage should narrow, and eventually reverse, as the state-action space grows toward Chengdu's scale, where a lookup table's implicit reliance on each state-action pair being visited often enough to be estimated reliably becomes harder to satisfy, and where a GNN's ability to generalize across structurally similar but never-exactly-visited candidate regions should start to pay for itself. Testing that hypothesis is the single most important remaining experiment in this research programme, not a minor extension of it.

\textbf{On the third gap}, transfer across cities, this paper reports no result at all yet, honestly: the cross-city curriculum (Section~\ref{sec:method-curriculum}) requires a second city's training run to transfer \emph{from}, and Xi'an -- trained here from a random initialization, as the first city in the planned sequence -- has not yet been used to warm-start a second city's policy. We view this, along with the Chengdu-scale test above, as the paper's clearest and most concrete next step, not a gap papered over with speculation: the code, data pipeline, and training infrastructure for at least Ho Chi Minh City (Section~\ref{sec:data}) already exist and are tested; what remains is running it.

\textbf{Limitations.} Beyond the incomplete city coverage already discussed, four limitations of this result deserve explicit statement rather than being left implicit. First, our three method variants are each reported over three seeds against the tabular baseline's five, and three-seed standard deviations on demand are large enough (Table~\ref{tab:results}) that most between-variant differences on that metric are not statistically distinguishable at this sample size; a fuller seed sweep, matching or exceeding the tabular baseline's five, is a direct, low-cost next step (Table~\ref{tab:weight-sensitivity}'s single-seed weight-conditioning check is a qualitative existence proof -- ``the mechanism responds to $\mathbf{w}$ at all'' -- and is far less sensitive to this concern than the quantitative demand/equity/coverage comparison in Table~\ref{tab:results}). Second, our accessibility metric is a Chebyshev-distance coverage proxy, not \citeauthor{zhang2024}'s radiation-model measure; the two are related in spirit (both reward geographic spread of station placement) but not numerically comparable, and implementing the exact radiation model is near-term follow-up work rather than a fundamental obstacle. Third, the MetroGNN-like and Zhang-et-al-like baseline rows are principled ablations of our own architecture, not reproductions of either paper's own reported numbers on their own data and code; the comparison this licenses is about the marginal value of each design choice within our method, not a head-to-head claim that our reimplementation of \citeauthor{su2024}'s idea outperforms their published system. Fourth, this is a single-city result on Xi'an specifically, whose 29$\times$29 grid and 8-directional action-space convention (Section~\ref{sec:method-implementation}) were adopted to match \citeauthor{michailidis2026}'s own benchmark exactly; generalization to the more irregular, polygon-based candidate-region graphs of Hanoi and Ho Chi Minh City (Section~\ref{sec:data}) uses a structurally different but formally equivalent action-space realization that has not yet been exercised end-to-end.

None of these limitations, we think, undermine the paper's central methodological contributions -- the architecture, the weight-conditioning mechanism and the bug it fixes, the two concrete implementation pitfalls documented in Section~\ref{sec:experiments} -- which stand independently of how many cities and seeds have been run so far. They do mean that the scale-dependent and transfer-learning claims motivating this research programme (Section~\ref{sec1}) remain hypotheses this paper has argued for and instrumented cleanly, rather than results this paper has already established.

\section{Conclusion}\label{sec13}

Metro network expansion is a decision cities get to make only rarely and can essentially never undo, which makes the case for bringing more rigorous, quantitative methods to it -- and for being honest about exactly what those methods have and have not yet demonstrated -- unusually strong. This paper proposed a GNN-based, multi-objective actor-critic policy for the Metro Network Expansion Problem that combines spatial representation learning \citep{su2024}, fairness-aware multi-objective reward decomposition \citep{zhang2024}, and a direct, architectural response to the training-efficiency critique raised by \citet{michailidis2026}, together with a rail-only scope discipline that lets the method extend, for the first time in this literature, to small and data-sparse emerging metro systems by treating each city's existing bus network as an auxiliary input signal rather than an additional design target.

We reported a first, fully real, GPU-trained experimental result on Xi'an: our full method outperforms both of its own ablations on every metric measured, and a direct behavioral sanity check confirms the weight-conditioning mechanism that makes ``one policy, many trade-offs'' true rather than aspirational -- but only after we caught and fixed a bug that had silently made it false in an earlier, otherwise-unremarkable-looking training run, a finding we report in detail because we believe it generalizes to other implementations of weight-conditioned multi-objective RL. On the narrower question of raw demand satisfaction under a matched, modest training budget, the tabular RL baseline of \citet{michailidis2026} outperforms our method on Xi'an, a result we report plainly rather than obscure, and read as consistent with rather than contradicting their central claim at this network's scale -- which is exactly why the scale-dependent comparison at Chengdu, and the cross-city curriculum linking Xi'an's trained weights to a second city, are the two most important next steps in this research programme, not afterthoughts to it. The infrastructure for both -- data pipelines, GNN encoder, multi-objective policy, and training loop -- already exists, is tested, and is described in full in this paper and its accompanying repository; what remains is running it, and we intend to.

\backmatter

\bmhead{Acknowledgements}

%% TODO: add acknowledgements if applicable.

\section*{Declarations}

\begin{itemize}
\item \textbf{Funding}: Not applicable. (Publication costs for \textit{Railway Engineering Science} are covered by Southwest Jiaotong University; no external funding has been received for this work.)
\item \textbf{Competing interests}: The author declares no competing interests.
\item \textbf{Ethics approval and consent to participate}: Not applicable -- this study uses publicly available transit network, demographic, and demand-proxy data and does not involve human or animal subjects.
\item \textbf{Consent for publication}: Not applicable.
\item \textbf{Data availability}: %% TODO: finalize once the data pipeline (Section~\ref{sec:data}) is implemented; link derived/preprocessed datasets in the project repository where redistribution licenses permit.
\item \textbf{Materials availability}: Not applicable.
\item \textbf{Code availability}: Code will be made available at \url{https://github.com/vinhqdang/metro_expansion} upon publication.
\item \textbf{Author contribution}: Q.-V.D. conceived the study, developed the method, and wrote the manuscript.
\end{itemize}

\begin{appendices}

\section{Extended hyperparameters and training details}\label{secA1}

\subsection{Compute environment}

All experiments were run in a conda environment (Python 3.13) with a CUDA-enabled PyTorch build (2.11, CUDA 12.8) and PyTorch Geometric (2.7.0), on a single NVIDIA RTX A2000 (8\,GB) laptop GPU. Geospatial preprocessing (Section~\ref{sec:data}) additionally uses geopandas (1.1.4), osmnx (2.1.1), and rasterio (1.5.1). Full dependency versions are pinned in \texttt{requirements.txt} in the project repository.

\subsection{Our method (full and ablations)}

\begin{table}[h]
\caption{Hyperparameters for the full method and the two ablations (Section~\ref{sec:experiments}), \texttt{scripts/train\_xian.py}}\label{tab:hparams-ours}
\begin{tabular}{@{}ll@{}}
\toprule
Hyperparameter & Value \\
\midrule
GNN hidden dimension & 32 \\
GNN message-passing layers ($L$) & 2 \\
Weight-embedding dimension & 32 (matches GNN hidden dimension) \\
Optimizer & Adam \\
Learning rate & $3 \times 10^{-4}$ \\
Gradient clipping (max norm) & 1.0 \\
Advantage normalization & per-episode (mean/std), when std $> 10^{-6}$ \\
Value-loss weight & 0.5 \\
Discount factor & 1.0 (undiscounted; finite-horizon episodic, Section~\ref{sec:problem}) \\
Tchebycheff $\varepsilon$ & $10^{-6}$ \\
Tchebycheff weight sampling & Dirichlet(1,1,1) each episode, i.e. uniform over the simplex \\
Training episodes & 500 (matched to the tabular RL baseline's budget) \\
Stations per episode ($nr\_stations$) & 20 \\
Price-equity groups ($nr\_groups$) & 5 \\
Coverage radius $K$ (Chebyshev) & 2 grid cells \\
Random seeds & 1, 2, 42 \\
\botrule
\end{tabular}
\end{table}

The \texttt{single\_objective} ablation is identical except the Tchebycheff weight vector is fixed to $(1,0,0)$ (demand only) every episode, rather than sampled. The \texttt{flat\_encoder} ablation is identical except the GNN encoder is replaced by \texttt{FlatEncoder} (\texttt{src/gnn/encoder.py}): a single linear projection of raw node features with $\tanh$ activation and no message passing over either edge type, isolating the contribution of graph structure from everything else held fixed.

\subsection{Tabular RL baseline (Michailidis et al.)}

\begin{table}[h]
\caption{Hyperparameters used for the tabular RL baseline (Section~\ref{sec:experiments}), \texttt{train\_qlearning.py} defaults from \citet{michailidis2026}'s public code release}\label{tab:hparams-tabular}
\begin{tabular}{@{}ll@{}}
\toprule
Hyperparameter & Value \\
\midrule
Learning rate $\alpha$ & 0.4 \\
Discount factor $\gamma$ & 0.8 \\
Exploration & $\varepsilon$-greedy, initial $\varepsilon = 1.0$, final $\varepsilon = 0.0$ \\
$\varepsilon$ decay steps & 400 \\
Update method & Temporal-difference (TD) \\
Reward type & \texttt{max\_efficiency} (demand only) \\
OD type & \texttt{pct} (percentage of total OD demand satisfied) \\
Training episodes & 500 \\
Test episodes & 10 \\
Stations per episode & 20 \\
Existing lines & included (not ignored) \\
Random seeds & 1, 2, 3, 7, 42 \\
\botrule
\end{tabular}
\end{table}

We use the baseline's own default hyperparameters throughout rather than tuning them in our favor or against it; the only values we set explicitly are \texttt{nr\_stations} and \texttt{train\_episodes}, matched to our own method's budget for a fair training-cost comparison. \texttt{docs/baseline\_reproduction.md} in the project repository records the exact commands, environment setup, and a real upstream bug found and patched in the baseline's own evaluation code (Section~\ref{sec:experiments}) needed to reproduce these numbers.

\end{appendices}

%%===========================================================================================%%
%% If you are submitting to one of the Nature Portfolio journals, using the eJP submission   %%
%% system, please include the references within the manuscript file itself. You may do this  %%
%% by copying the reference list from your .bbl file, paste it into the main manuscript .tex %%
%% file, and delete the associated \verb+\bibliography+ commands.                            %%
%%===========================================================================================%%

\bibliography{references}% common bib file
%% if required, the content of .bbl file can be included here once bbl is generated
%%\input sn-article.bbl

\end{document}
